\subsection{Classification fidelity and information retained by member thresholds}

\subsubsection{Classification fidelity}

The affine trigger errors are small, yet they can still change a classification when a price lies between the affine and benchmark thresholds. The comparison is therefore made on the price axis, and two quantities are computed for each configuration from the affine and benchmark member-threshold vectors. The first is the intersection over union (IoU) of the affine and benchmark contested-price intervals. The second is the omitted mixed share, the length of the benchmark interval outside the affine interval divided by the benchmark width. For prices in this omitted region, the affine rule reports a unanimous classification, whereas the benchmark reports a mixed one. Both diagnostics are geometric, so they require no assumption on the distribution of the prevailing price.

Both diagnostics are evaluated on a fixed deterministic design, with the reference project value equal to the investment cost. The design combines the 33 admissible calibrations with four panel sizes (3, 5, 7, and 9) and six recorded-weight profiles. Crossed with relative radii of 5\%, 10\%, 15\%, and 22\%, these factors give 3,168 configurations.

Over this design, Criterion C1 requires at least 95\% of configurations to attain an IoU of at least 0.90. Criterion C2 requires the 95th percentile of the omitted mixed share to be at most 5\%. Both criteria therefore concern the tail of the distribution across configurations. On the full set, 89.5\% of configurations attained an IoU of at least 0.90, below the 95\% required by C1. The 95th percentile of the omitted mixed share was 7.65\%, above the C2 limit, so the full set met neither criterion. Fidelity improved at smaller radii, and at the 10\% radius alone, 98.2\% of configurations met the IoU requirement. On the pooled 5\% and 10\% set, this share was 99.1\%, and the 95th percentile of the omitted mixed share was 4.15\%. Both criteria held on this restricted set.

\subsubsection{Information retention}

Criteria C1 and C2 test whether the affine interval reproduces the benchmark interval. A separate simulation examines whether the affine member-threshold vector retains price-axis information lost by any single-trigger summary of the benchmark vector. The simulation generates 2,000 panels for each combination of five radii, four panel sizes, and five weight families, giving 200,000 panels. Of these, 160,000 lie within the primary radius range. Each simulated panel draws one of the 33 calibrations.

Member-vector loss is the mean absolute price-axis distance between each affine member threshold and its benchmark counterpart, expressed as a percentage of the anchor threshold. Four single-trigger candidates enter the comparison: the anchor threshold, the benchmark thresholds at the mean and median recorded weights, and the largest benchmark member threshold. The loss of a candidate is computed by replacing every affine member threshold with the candidate threshold. For each panel, a four-candidate ex post comparator uses the generated benchmark vector to select the candidate with the lowest loss. This hindsight choice favors the single-trigger summaries in the comparison.

Criterion C3 has two parts. The median member-vector loss must be at most one half of the median loss of the comparator. The four primary radii, four panel sizes, and five weight families define 80 design cells. Each cell contributes the median, over its panels, of the comparator loss minus the member-vector loss. The 95\% resampling interval of these cell-level differences must also exclude zero.

Across the panels within the primary radius range, the median member-vector loss was 0.224\% of the anchor threshold. The median loss of the comparator was 1.944\%, and the resulting 88.5\% reduction in median loss met the first part of C3. The member-threshold vector had the lower loss in 99.5\% of the panels.

The comparator selected the median-weight benchmark threshold in every tested panel, since the median minimizes an absolute-distance loss and all tested panel sizes are odd. For the second part of Criterion C3, the median cell-level loss difference across the 80 design cells was 1.65\% of the anchor threshold. Its 95\% resampling interval ran from 1.31\% to 2.06\% and excluded zero, so Criterion C3 was met.

A separate deterministic grid, equal in size to the fidelity design, examined how much of this gain requires the full vector. A three-threshold summary, reporting the smallest, median, and largest benchmark thresholds, cut the median loss to 0.502\% of the anchor threshold. On the same grid, the median loss was 2.240\% for the best single-trigger summary and 0.235\% for the member-threshold vector. The full vector thus removed 53\% of the loss left by the three-threshold summary and was more accurate in 72.5\% of configurations. Three-member panels form a quarter of this grid, and there the three-threshold summary reproduces the benchmark vector exactly. The comparison is therefore informative only on the remaining configurations. Adding a min-max midpoint and a quartile-based composite to the single-trigger candidates left this ordering of losses unchanged. Under both candidate sets, most of the improvement over a single-trigger summary came from reporting more than one threshold. The member-threshold vector also identifies the member behind each threshold, whereas a sorted summary loses this link.

For the Monte Carlo panels, Table~\ref{tab:simulation} reports the results by radius, and each row summarizes 40,000 simulated panels. Its omitted-mixed-share percentiles differed from the deterministic statistics because the simulation drew random weight profiles.

The same panels supply two reference quantities for an institution setting its own release criteria. The first is a length in price units, the omitted mixed share multiplied by the benchmark interval width. Table~\ref{tab:simulation} reports this length normalized by the anchor threshold, from 0.09\% at a 5\% radius to 3.60\% at the stress radius. At the baseline anchor threshold, these values equaled about 0.10 and 4.24\,CNY/t. The second is panel dependence. Fidelity improved monotonically with panel size and degraded under one-sided weight skew. A skewed profile brings the two interval edges closer together, so a given threshold error removes a larger share of the benchmark interval. The weakest cell was a three-member panel under skew. A committee can compare these lengths with the price band over which it regards the gate decision as open. Given the panel dependence, it can then set its own release cap.

The weakest configuration, a three-member panel under skew, is visible at the gate without the benchmark. A dispersion statistic of the recorded weights might therefore serve as an escalation rule for benchmark recomputation. One natural candidate is the ratio of the range of the recorded weights to their interquartile range. This ratio failed as an escalation rule. For three-member panels, it equals two for every profile. Their classification fidelity nonetheless ranged from 48.5\% to 99.2\%, measured as the share of configurations with an IoU of at least 0.90. For larger panels, the correlation between the ratio and fidelity stayed below 0.25. The weight profile and the panel size did separate passing cells from failing cells, and a gate records both. The release rule nonetheless conditions on the distances between the prevailing price and the thresholds, because these distances catch the affected gates directly.

Beyond the simulated panels, the four-member reference profile was evaluated deterministically in 30 direct comparisons with the benchmark, over five radii and six reference values. These checks and the 32 archived-price cells of Table~\ref{tab:margins} are included in the replication package.

\begin{table*}[!t]
\caption{Classification fidelity and price-axis information loss across simulated panels of sizes 3, 5, 7, and 9. The 33\% radius is a stress case.}
\label{tab:simulation}
\centering
\footnotesize
\begin{tabular}{@{}lL{0.17\textwidth}L{0.3\textwidth}L{0.33\textwidth}@{}}
\toprule
$h/\sbar$ & IoU, median (5th percentile) & 95th-percentile omitted mixed share / anchor-normalized omitted length (\%) & Median loss, member vector / four-candidate ex post comparator (\% of anchor threshold)\\
\midrule
5\% & 0.978 (0.929) & 4.0 / 0.09 & 0.039 / 0.862\\
10\% & 0.957 (0.862) & 7.9 / 0.35 & 0.156 / 1.720\\
15\% & 0.937 (0.801) & 11.7 / 0.78 & 0.350 / 2.601\\
22\% & 0.911 (0.723) & 16.8 / 1.65 & 0.735 / 3.810\\
33\% & 0.876 (0.621) & 23.9 / 3.60 & 1.573 / 5.746\\
\bottomrule
\end{tabular}
\end{table*}

\subsubsection{Finite-registry safeguard}

Within the tested design, the simulation supported retaining the member-threshold vector. Classification fidelity near the interval edges, however, weakened as the radius grew, most sharply for skewed profiles and small panels. A release rule is therefore attached to the member-threshold report. Consider a registered calibration $\theta=(r,\mu_c,\sbar)$ and an evaluated radius $\rho=h/\sbar$. With $L_{\mathrm{lin}}$ and $U_{\mathrm{lin}}$ as the edges of the affine interval $C_t$, let
\begin{equation}
\begin{aligned}
\eR(\theta,\rho)&:=\max_{\alpha\in[0,1]}\frac{\lvert\Vlin(\alpha)-\VMEU(\alpha,I)\rvert}{\Vzero},\\
\ER(\theta,\rho)&:=P^{*}_{0}\,[\eR(\theta,\rho)+\varepsilon_{\mathrm{num}}],\\
\dedge&:=\min\{\lvert P_t-L_{\mathrm{lin}}\rvert,\;\lvert P_t-U_{\mathrm{lin}}\rvert\},\\
\dmember&:=\min_i\,\lvert P_t-P^{*}_{i,\mathrm{lin}}\rvert.
\end{aligned}
\label{eq:safeguard}
\end{equation}
The distance $\dedge$ protects the three-state gate summary of unanimous wait, mixed, or unanimous invest, whereas $\dmember$ protects the invest-or-wait classification of every member. The interval edges are member thresholds, so $\dmember\le\dedge$ and the member condition binds, although both distances are retained for audit. Each distance is compared with the envelope $\ER$, which converts the registered maximum anchor-normalized member-trigger error plus a numerical tolerance into price units. Because $P^{*}_{0}$ multiplies both $\eR$ and the tolerance, a common rescaling of the member thresholds rescales $\ER$ by the same factor. The anchor-normalized trigger error is invariant, yet the release decision can still change, because the prevailing price stays fixed while the thresholds rescale. The registry of evaluated configurations therefore stores the dimensionless envelope $\eR$. After any revision of the investment cost or the captured volume, the gate recomputes $P^{*}_{0}$ and $\ER$ and re-evaluates both proximity distances. The audit record carries $I/Q$ and $P^{*}_{0}$ with the registry identifiers, so a past release can be reconstructed on its original price axis.

The registry decides only when the benchmark may be omitted, and a committee may still request the benchmark at any gate. Affine-only reporting requires four conditions. (i) The calibration and the canonical radius key are registered exactly. (ii) The radius is at most 10\%. (iii) The reference value satisfies $V_0/I=1$, the only value at which the release rule was evaluated. (iv) Both proximity distances exceed $\ER$. If any condition fails, or if the solver or envelope version is unregistered, the gate recomputes the benchmark and issues a dual report. In either mode, the reporting interface passes on the member-threshold vector and the fallback status. A registry match under condition (i) must come from the frozen serialized key. Rounded display values and values reparsed from the upstream signal file do not qualify. Interpolation is excluded, because no certified release envelope exists between evaluated radius keys. The 10\% radius cap is a hard cutoff, so rounding a configuration constant can move a radius to the wrong side of it. One low-radius scenario targeted exactly 10\% and yielded 10.30\% at its rounded constant, outside the affine-only region. Exact-key matching guards against this error. An institution may prespecify a rule mapping each gate radius to a registered key, and the release guarantee then applies to the discretized ambiguity set. The audit record retains both the original radius and the registered radius.

When the benchmark is recomputed, the gate first evaluates the uniqueness condition~(\ref{eq:unique}). When this sufficient condition holds, the maximizer is the unique stationary point on $[\Vlow,\Vhigh]$, and a bounded solve suffices. When it fails, the gate isolates every stationary point and compares the objective at each and at both endpoints before releasing a benchmark classification. Condition~(\ref{eq:onesided}) is recorded separately, because it governs the one-sided reading of the dual report. Algorithm~\ref{alg:gate} states the matching, versioning, and solver rules.

The registry makes a released classification auditable after the fact. Each registry cell is indexed by the economic calibration, the radius key, and the reference value, and it carries the envelope and solver versions. The gate audit record also holds the recorded weights and the prevailing price. A later reviewer can then identify the inputs behind a past classification and recompute it without re-eliciting any weight. In prospective use, the registry should be append-only, so a revised input opens a new dated cell. The earlier cell then stays readable with its original guarantee. The release rule also isolated the configurations in which the two formulations disagreed. Across the 3,960 registered configurations, such disagreement occurred in 57. None of these 57 passed the radius cap and both proximity conditions together. Two of them lay at a 10\% radius, where a proximity flag routed them to benchmark recomputation. The envelopes behind these flags were constructed at all 3,960 domain-valid configurations with $V_0/I=1$, and two checks bounded their numerical error. First, a denser $\alpha$ grid detected no exceedance of the registered envelope in any of the 165 calibration-radius cells. Second, the maximum on the finite weight grid was compared with a certified upper bound on the discrepancy over $[0,1]$ in each cell. This bound used a closed-form Lipschitz constant and a branch-and-bound search. It exceeded the grid maximum by at most $3.1\times10^{-6}$\,CNY/t and changed no release decision.

The 10\% radius cap in condition (ii) is the largest evaluated radius at which both fidelity criteria held in the deterministic design. A refinement at intermediate radii located the failure just above the cap. C2 failed first, at 11\%, where the 95th percentile of the omitted mixed share reached 5.4\% against the 5\% limit. At the 10\% radius, the same quantity was 4.96\%, so C2 held at the cap by 0.04 percentage points. The breach at 11\% was 0.4 percentage points, which suggests a gradual loss of accuracy across the cap. By Proposition~\ref{prop:accuracy}(i), the leading error term grows with the squared radius, so raising the cap from 10\% to 22\% multiplies it by about five.

The cap was also tested on 18 calibrations and five weight profiles absent from the evaluated grid, six of these calibrations lying outside its range. The price-axis quantity used for release behaved similarly on this held-out set. The 95th percentile across cells of the largest within-cell member-threshold deviation was 0.5141\,CNY/t, against 0.5143\,CNY/t on the evaluated grid. The normalized fidelity criteria were more sensitive to the weight profile and held on four of the five held-out profiles. They failed on one profile with weights clustered within 0.06 of one half, for which the benchmark interval shrank to 0.72\% of the anchor threshold. A similar absolute discrepancy then became a much larger relative one. Both criteria registered this effect, because IoU normalizes by the union of the two intervals and the omitted mixed share by the benchmark width. For this reason, the release rule conditions on absolute price-axis distances.

Above the cap, the dual report is mandatory. Across the three radii above the cap and the weakest profiles, member-level classifications agreed in more than nine cases in ten. The remaining differences ran in one direction. In no evaluated cell did the affine rule classify a price as invest while the benchmark classified it as wait. Proposition~\ref{prop:benchmark}(ii) explains this pattern, because under the one-sidedness condition the affine report errs toward deferral. The condition held throughout the affine-only region as well, so any residual affine error there also leans toward deferral.

Figure~\ref{fig:release} shows the release decision as a single pass over the four release conditions. When all four hold, $\ER$ bounds the error of every member threshold, so the released classification matches the benchmark classification. The released interval width $B_t$, however, is bounded only in absolute terms, within $2\ER$. In relative terms, the width error for the four-member reference profile was 1.29\% at the reference project value of the release rule. The 2.25\% maximum across tested reference values arose outside the release region. The width shrinks as the recorded weights cluster, so the same absolute error is a larger fraction of the width on tightly clustered profiles. In Figure~\ref{fig:release}, input validity is checked on a separate branch, and a failed upstream acceptance check blocks the threshold report.

With the release rule in place, the model results fall into two groups with different scopes. The analytical group, covering the critical weight and its location relative to one half, holds wherever the maintained model holds. The numerical group, covering the accuracy of the affine rule and its release conditions, holds only on the finite registry of evaluated configurations. The practical relevance of both groups depends on the inputs of a specific project.

\begin{figure*}[!t]
\centering
\resizebox{0.97\textwidth}{!}{%
\begin{tikzpicture}[
  font=\footnotesize,
  proc/.style={draw, rounded corners=2pt, align=center, fill=blue!7, minimum height=0.9cm},
  dec/.style={draw, diamond, aspect=2.1, align=center, fill=orange!10, inner sep=1pt},
  ok/.style={draw, rounded corners=2pt, align=center, fill=teal!12, minimum height=0.9cm},
  arr/.style={-{Stealth[length=2mm]}, thick},
  lab/.style={font=\scriptsize\itshape}
]
\node[proc, text width=2.0cm] (in) at (0,0) {Inputs pass\\acceptance checks};
\node[dec] (reg) at (3.2,0) {Keys registered\\exactly};
\node[dec] (rad) at (6.2,0) {$\rho\le10\%$};
\node[dec] (v0) at (9.0,0) {$V_0/I=1$};
\node[dec] (dist) at (12.4,0) {$\dedge,\dmember$\\both $>\ER$};
\node[ok, text width=3.0cm] (aff) at (16.3,0) {Affine-only release of\\member thresholds,\\$C_t$, and $B_t$};
\node[proc, text width=2.6cm] (rev) at (0,-2.6) {Input review required;\\no threshold report};
\node[proc, text width=4.2cm] (rec) at (6.6,-2.6) {Recompute the $\aMEU$ benchmark\\at every recorded weight};
\node[proc, text width=4.2cm] (dual) at (12.8,-2.6) {Dual report\\summaries agree $\Rightarrow$ release\\the common classification;\\summaries differ $\Rightarrow$\\benchmark-model review,\\no classification released};
\draw[arr] (in) -- (reg);
\draw[arr] (reg) -- node[lab, above]{yes} (rad);
\draw[arr] (rad) -- node[lab, above]{yes} (v0);
\draw[arr] (v0) -- node[lab, above]{yes} (dist);
\draw[arr] (dist) -- node[lab, above]{yes} (aff);
\draw[arr] (in) -- node[lab, left]{fail} (rev);
\draw[arr] (reg.south) -- node[lab, left, pos=0.3]{no} (rec.north west);
\draw[arr] (rad.south) -- node[lab, left, pos=0.4]{no} (rec.north);
\draw[arr] (v0.south) -- node[lab, right, pos=0.3]{no} (rec.north east);
\draw[arr] (dist.south) -- node[lab, right, pos=0.3]{no} ($(rec.north east)+(0.1,-0.25)$);
\draw[arr] (rec) -- (dual);
\node[lab, align=left, anchor=north west] at ($(dual.south west)+(0,-0.15)$) {Directional labels are released with $V_0$ and de-emphasized\\inside the review band around $\tilde{\alpha}$.};
\end{tikzpicture}}
\caption{Release decision for one gate. The gate first runs the registered input acceptance checks, and a failure blocks the threshold report. Four release conditions are then checked in sequence. The calibration and radius key must be registered exactly, and the relative radius must be at most 10\%. The reference project value must satisfy $V_0/I=1$, and both $\dedge$ and $\dmember$ must exceed the envelope $\ER$. Passing all four conditions releases the affine-only report with the member-threshold vector, $C_t$, and $B_t$. Failing any condition routes the gate to a dual report, with the $\aMEU$ benchmark recomputed at every recorded weight. In this mode, both member-threshold vectors are issued. A gate-level classification is released only when the two three-state summaries agree, although individual member classifications may still differ. A discrepancy between the summaries leads to a benchmark-model review, with no precedence for either rule.}
\label{fig:release}
\end{figure*}
